% Options for packages loaded elsewhere
% Options for packages loaded elsewhere
\PassOptionsToPackage{unicode}{hyperref}
\PassOptionsToPackage{hyphens}{url}
\PassOptionsToPackage{dvipsnames,svgnames,x11names}{xcolor}
%
\documentclass[
  brazilian,
  letterpaper,
  DIV=11,
  numbers=noendperiod]{scrreprt}
\usepackage{xcolor}
\usepackage[margin=2.5cm]{geometry}
\usepackage{amsmath,amssymb}
\setcounter{secnumdepth}{5}
\usepackage{iftex}
\ifPDFTeX
  \usepackage[T1]{fontenc}
  \usepackage[utf8]{inputenc}
  \usepackage{textcomp} % provide euro and other symbols
\else % if luatex or xetex
  \usepackage{unicode-math} % this also loads fontspec
  \defaultfontfeatures{Scale=MatchLowercase}
  \defaultfontfeatures[\rmfamily]{Ligatures=TeX,Scale=1}
\fi
\usepackage{lmodern}
\ifPDFTeX\else
  % xetex/luatex font selection
\fi
% Use upquote if available, for straight quotes in verbatim environments
\IfFileExists{upquote.sty}{\usepackage{upquote}}{}
\IfFileExists{microtype.sty}{% use microtype if available
  \usepackage[]{microtype}
  \UseMicrotypeSet[protrusion]{basicmath} % disable protrusion for tt fonts
}{}
\makeatletter
\@ifundefined{KOMAClassName}{% if non-KOMA class
  \IfFileExists{parskip.sty}{%
    \usepackage{parskip}
  }{% else
    \setlength{\parindent}{0pt}
    \setlength{\parskip}{6pt plus 2pt minus 1pt}}
}{% if KOMA class
  \KOMAoptions{parskip=half}}
\makeatother
% Make \paragraph and \subparagraph free-standing
\makeatletter
\ifx\paragraph\undefined\else
  \let\oldparagraph\paragraph
  \renewcommand{\paragraph}{
    \@ifstar
      \xxxParagraphStar
      \xxxParagraphNoStar
  }
  \newcommand{\xxxParagraphStar}[1]{\oldparagraph*{#1}\mbox{}}
  \newcommand{\xxxParagraphNoStar}[1]{\oldparagraph{#1}\mbox{}}
\fi
\ifx\subparagraph\undefined\else
  \let\oldsubparagraph\subparagraph
  \renewcommand{\subparagraph}{
    \@ifstar
      \xxxSubParagraphStar
      \xxxSubParagraphNoStar
  }
  \newcommand{\xxxSubParagraphStar}[1]{\oldsubparagraph*{#1}\mbox{}}
  \newcommand{\xxxSubParagraphNoStar}[1]{\oldsubparagraph{#1}\mbox{}}
\fi
\makeatother

\usepackage{color}
\usepackage{fancyvrb}
\newcommand{\VerbBar}{|}
\newcommand{\VERB}{\Verb[commandchars=\\\{\}]}
\DefineVerbatimEnvironment{Highlighting}{Verbatim}{commandchars=\\\{\}}
% Add ',fontsize=\small' for more characters per line
\usepackage{framed}
\definecolor{shadecolor}{RGB}{241,243,245}
\newenvironment{Shaded}{\begin{snugshade}}{\end{snugshade}}
\newcommand{\AlertTok}[1]{\textcolor[rgb]{0.68,0.00,0.00}{#1}}
\newcommand{\AnnotationTok}[1]{\textcolor[rgb]{0.37,0.37,0.37}{#1}}
\newcommand{\AttributeTok}[1]{\textcolor[rgb]{0.40,0.45,0.13}{#1}}
\newcommand{\BaseNTok}[1]{\textcolor[rgb]{0.68,0.00,0.00}{#1}}
\newcommand{\BuiltInTok}[1]{\textcolor[rgb]{0.00,0.23,0.31}{#1}}
\newcommand{\CharTok}[1]{\textcolor[rgb]{0.13,0.47,0.30}{#1}}
\newcommand{\CommentTok}[1]{\textcolor[rgb]{0.37,0.37,0.37}{#1}}
\newcommand{\CommentVarTok}[1]{\textcolor[rgb]{0.37,0.37,0.37}{\textit{#1}}}
\newcommand{\ConstantTok}[1]{\textcolor[rgb]{0.56,0.35,0.01}{#1}}
\newcommand{\ControlFlowTok}[1]{\textcolor[rgb]{0.00,0.23,0.31}{\textbf{#1}}}
\newcommand{\DataTypeTok}[1]{\textcolor[rgb]{0.68,0.00,0.00}{#1}}
\newcommand{\DecValTok}[1]{\textcolor[rgb]{0.68,0.00,0.00}{#1}}
\newcommand{\DocumentationTok}[1]{\textcolor[rgb]{0.37,0.37,0.37}{\textit{#1}}}
\newcommand{\ErrorTok}[1]{\textcolor[rgb]{0.68,0.00,0.00}{#1}}
\newcommand{\ExtensionTok}[1]{\textcolor[rgb]{0.00,0.23,0.31}{#1}}
\newcommand{\FloatTok}[1]{\textcolor[rgb]{0.68,0.00,0.00}{#1}}
\newcommand{\FunctionTok}[1]{\textcolor[rgb]{0.28,0.35,0.67}{#1}}
\newcommand{\ImportTok}[1]{\textcolor[rgb]{0.00,0.46,0.62}{#1}}
\newcommand{\InformationTok}[1]{\textcolor[rgb]{0.37,0.37,0.37}{#1}}
\newcommand{\KeywordTok}[1]{\textcolor[rgb]{0.00,0.23,0.31}{\textbf{#1}}}
\newcommand{\NormalTok}[1]{\textcolor[rgb]{0.00,0.23,0.31}{#1}}
\newcommand{\OperatorTok}[1]{\textcolor[rgb]{0.37,0.37,0.37}{#1}}
\newcommand{\OtherTok}[1]{\textcolor[rgb]{0.00,0.23,0.31}{#1}}
\newcommand{\PreprocessorTok}[1]{\textcolor[rgb]{0.68,0.00,0.00}{#1}}
\newcommand{\RegionMarkerTok}[1]{\textcolor[rgb]{0.00,0.23,0.31}{#1}}
\newcommand{\SpecialCharTok}[1]{\textcolor[rgb]{0.37,0.37,0.37}{#1}}
\newcommand{\SpecialStringTok}[1]{\textcolor[rgb]{0.13,0.47,0.30}{#1}}
\newcommand{\StringTok}[1]{\textcolor[rgb]{0.13,0.47,0.30}{#1}}
\newcommand{\VariableTok}[1]{\textcolor[rgb]{0.07,0.07,0.07}{#1}}
\newcommand{\VerbatimStringTok}[1]{\textcolor[rgb]{0.13,0.47,0.30}{#1}}
\newcommand{\WarningTok}[1]{\textcolor[rgb]{0.37,0.37,0.37}{\textit{#1}}}

\usepackage{longtable,booktabs,array}
\usepackage{calc} % for calculating minipage widths
% Correct order of tables after \paragraph or \subparagraph
\usepackage{etoolbox}
\makeatletter
\patchcmd\longtable{\par}{\if@noskipsec\mbox{}\fi\par}{}{}
\makeatother
% Allow footnotes in longtable head/foot
\IfFileExists{footnotehyper.sty}{\usepackage{footnotehyper}}{\usepackage{footnote}}
\makesavenoteenv{longtable}
\usepackage{graphicx}
\makeatletter
\newsavebox\pandoc@box
\newcommand*\pandocbounded[1]{% scales image to fit in text height/width
  \sbox\pandoc@box{#1}%
  \Gscale@div\@tempa{\textheight}{\dimexpr\ht\pandoc@box+\dp\pandoc@box\relax}%
  \Gscale@div\@tempb{\linewidth}{\wd\pandoc@box}%
  \ifdim\@tempb\p@<\@tempa\p@\let\@tempa\@tempb\fi% select the smaller of both
  \ifdim\@tempa\p@<\p@\scalebox{\@tempa}{\usebox\pandoc@box}%
  \else\usebox{\pandoc@box}%
  \fi%
}
% Set default figure placement to htbp
\def\fps@figure{htbp}
\makeatother


% definitions for citeproc citations
\NewDocumentCommand\citeproctext{}{}
\NewDocumentCommand\citeproc{mm}{%
  \begingroup\def\citeproctext{#2}\cite{#1}\endgroup}
\makeatletter
 % allow citations to break across lines
 \let\@cite@ofmt\@firstofone
 % avoid brackets around text for \cite:
 \def\@biblabel#1{}
 \def\@cite#1#2{{#1\if@tempswa , #2\fi}}
\makeatother
\newlength{\cslhangindent}
\setlength{\cslhangindent}{1.5em}
\newlength{\csllabelwidth}
\setlength{\csllabelwidth}{3em}
\newenvironment{CSLReferences}[2] % #1 hanging-indent, #2 entry-spacing
 {\begin{list}{}{%
  \setlength{\itemindent}{0pt}
  \setlength{\leftmargin}{0pt}
  \setlength{\parsep}{0pt}
  % turn on hanging indent if param 1 is 1
  \ifodd #1
   \setlength{\leftmargin}{\cslhangindent}
   \setlength{\itemindent}{-1\cslhangindent}
  \fi
  % set entry spacing
  \setlength{\itemsep}{#2\baselineskip}}}
 {\end{list}}
\usepackage{calc}
\newcommand{\CSLBlock}[1]{\hfill\break\parbox[t]{\linewidth}{\strut\ignorespaces#1\strut}}
\newcommand{\CSLLeftMargin}[1]{\parbox[t]{\csllabelwidth}{\strut#1\strut}}
\newcommand{\CSLRightInline}[1]{\parbox[t]{\linewidth - \csllabelwidth}{\strut#1\strut}}
\newcommand{\CSLIndent}[1]{\hspace{\cslhangindent}#1}

\ifLuaTeX
\usepackage[bidi=basic]{babel}
\else
\usepackage[bidi=default]{babel}
\fi
% get rid of language-specific shorthands (see #6817):
\let\LanguageShortHands\languageshorthands
\def\languageshorthands#1{}


\setlength{\emergencystretch}{3em} % prevent overfull lines

\providecommand{\tightlist}{%
  \setlength{\itemsep}{0pt}\setlength{\parskip}{0pt}}



 


\KOMAoption{captions}{tableheading}
\makeatletter
\@ifpackageloaded{bookmark}{}{\usepackage{bookmark}}
\makeatother
\makeatletter
\@ifpackageloaded{caption}{}{\usepackage{caption}}
\AtBeginDocument{%
\ifdefined\contentsname
  \renewcommand*\contentsname{Índice}
\else
  \newcommand\contentsname{Índice}
\fi
\ifdefined\listfigurename
  \renewcommand*\listfigurename{Lista de Figuras}
\else
  \newcommand\listfigurename{Lista de Figuras}
\fi
\ifdefined\listtablename
  \renewcommand*\listtablename{Lista de Tabelas}
\else
  \newcommand\listtablename{Lista de Tabelas}
\fi
\ifdefined\figurename
  \renewcommand*\figurename{Figura}
\else
  \newcommand\figurename{Figura}
\fi
\ifdefined\tablename
  \renewcommand*\tablename{Tabela}
\else
  \newcommand\tablename{Tabela}
\fi
}
\@ifpackageloaded{float}{}{\usepackage{float}}
\floatstyle{ruled}
\@ifundefined{c@chapter}{\newfloat{codelisting}{h}{lop}}{\newfloat{codelisting}{h}{lop}[chapter]}
\floatname{codelisting}{Listagem}
\newcommand*\listoflistings{\listof{codelisting}{Lista de Listagens}}
\makeatother
\makeatletter
\makeatother
\makeatletter
\@ifpackageloaded{caption}{}{\usepackage{caption}}
\@ifpackageloaded{subcaption}{}{\usepackage{subcaption}}
\makeatother
\usepackage{bookmark}
\IfFileExists{xurl.sty}{\usepackage{xurl}}{} % add URL line breaks if available
\urlstyle{same}
\hypersetup{
  pdftitle={Introdução ao Python e Web Scraping},
  pdfauthor={Gabriel Sanábio Mazzetti Affonso},
  pdflang={pt-BR},
  colorlinks=true,
  linkcolor={blue},
  filecolor={Maroon},
  citecolor={Blue},
  urlcolor={Blue},
  pdfcreator={LaTeX via pandoc}}


\title{Introdução ao Python e Web Scraping}
\author{Gabriel Sanábio Mazzetti Affonso}
\date{2026-09-05}
\begin{document}
\maketitle

\renewcommand*\contentsname{Índice}
{
\hypersetup{linkcolor=}
\setcounter{tocdepth}{2}
\tableofcontents
}

\bookmarksetup{startatroot}

\chapter*{Prefácio}\label{prefuxe1cio}
\addcontentsline{toc}{chapter}{Prefácio}

\markboth{Prefácio}{Prefácio}

O presente material consiste na consolidação dos resultados de pesquisa
obtidos durante o período de Iniciação Científica, no contexto do
projeto \textbf{Métodos de Aprendizado Estatístico para Discriminação e
Classificação: Aplicações à Indústria 4.0 -- Fase 3}. Este documento
visa não só o detalhamento dos métodos de coleta de dados via \emph{Web
Scraping}, mas também a viabilidade do uso do \emph{Python} para tal,
além de destacar as especificidades da linguagem e suas vantagens para o
aprendizado de métodos estatísticos e computacionais.

Neste contexto, vale ressaltar que o material é de cunho introdutório,
buscando viabilizar o aprendizado do \emph{Python} enquanto linguagem de
programação e do \emph{Web Scraping} como técnica poderosa de coleta de
dados. Dessa forma, seu foco é desenvolver, no estatístico, o interesse
pelos diversos recursos computacionais que podem auxiliá-los enquanto
profissionais e estudiosos da área..

A motivação para o desenvolvimento deste material surgiu da necessidade
de diversificar as ferramentas de aplicação computacional existentes
para analisar dados. A interseção entre a teoria estatística e a
aplicação dos métodos se sustenta no casamento entre o rigor matemático
e a modernidade computacional. Qualquer das partes que for negligenciada
ocasionará problemas ou limitações importantes.

Ademais, é válido ressaltar a seção extra presente no material,
desenvolvida com o intuito de explicar a estrutura dos livros digitais,
como esse, usando o \emph{Quarto}. Ainda que muitos considerem adequado
o uso do \emph{Bookdown} para este tipo de trabalho (o que é totalmente
válido, dadas as especificidades do Rmarkdown), o \emph{Quarto} se
mostrou uma espécie de sucessor natural muito adequado, apresentando
ferramentas modernas e extremamente úteis para detalhes de
\emph{layout}, especialmente tratando de documentos compilados em
\emph{HTML}.

Por fim, deixo aqui minha gratidão à Universidade Federal de Juiz de
Fora (UFJF) pelo ambiente acadêmico. Agradeço, especialmente, ao meu
orientador, Professor Doutor Lupércio França Bessegato, pela
oportunidade e pelo acompanhamento técnico que possibilitaram o
desenvolvimento deste trabalho.

\bookmarksetup{startatroot}

\chapter{Introdução}\label{introduuxe7uxe3o}

Na prática, a estatística contemporânea vai muito além do domínio
teórico dos fundamentos, da modelagem probabilística e da inferência. A
sociedade atual convive com uma realidade em que os dados são o produto
mais valioso existente, além simplesmente de fontes de informação. Dessa
forma, observa-se nos últimos anos o surgimento de conceitos como o
\emph{Big Data}, para descrever o fenômeno da quantidade massiva de
dados que alteram significamente o fluxo de trabalho do analista de
dados, exigindo deste profissional, muitas vezes, muito além da
capacidade de analisar dados. É necessário ter aptidão para manusear os
dados apropriadamente, buscá-los em diferentes fontes, criar filtros
especificos relativos às necessidades de um determinado problema, entre
outras habilidades.

Neste contexto, este material foi elaborado com o intuito de introduzir
a linguagem \emph{Python} no arsenal analítico de estudantes e
profissionais de Estatística, demonstrando sua viabilidade e potência
para a análise de dados e a automação de rotinas de trabalho.

\section{A importância do Python na
Estatística}\label{a-importuxe2ncia-do-python-na-estatuxedstica}

Tradicionalmente, a linguagem \emph{R} ocupa um lugar importante na
formação do estatístico, por ser uma linguagem vetorizada e já
desenvolvida para o trabalho com dados. Porém, apesar do poder do
\emph{R} e seus pacotes nos trabalhos de modelagem de dados e no uso de
testes clássicos (paramétricos e não-paramétricos), o \emph{Python}
consolidou-se como uma linguagem extremamente versátil, que permite
aplicações robustas (como as de \emph{Machine Learning}, por exemplo) e
a integração de sistemas.

Em termos de aprendizado, o \emph{Python} é facilmente assimilado por
quem já teve contato com algoritmos, por conta de sua similaridade com
linguagens compiladas, como no caso do \emph{C++} (exigindo apenas
adaptações sintatícas), e com linguagens interpretadas, como o \emph{R}.

Na prática, o \emph{Python} apresenta sua maior vantagem ao analisarmos
sua quantidade de usuários e, consequentemente, sua quantidade de
bibliotecas desenvolvidas para diferentes aplicações. No contexto da
Estatística, destacamos o \emph{pandas}, para estruturação de dados
tabulares, e o \emph{numpy}, para computação numérica de alta
performance. O domínio dessas ferramentas permite ao analista processar
grandes volumes de dados com eficiência, além de explorar diferentes
possibilidades, como a de romper com a dependência de bancos de dados
pré-estruturados através da automatização da coleta de dados via
\emph{internet}.

\section{O escopo deste material}\label{o-escopo-deste-material}

Este material tem um teor introdutório e prático. Ele não tem a intenção
de esgotar as funcionalidades do \emph{Python} (que seria, inclusive,
uma tarefa praticamente impossível), tampouco de aprofundar-se nas
teorias estatísticas que fundamentam nossas análises ou na teoria
computacional que fundamenta e explica as especificidades da linguagem.
O objetivo central é ser uma ponte rápida para profissionais de
Estatística compreenderem a linguagem e o ecossistema, e como podem se
beneficiar deles.

Neste âmbito, a abordagem adotada se resume em execução de códigos e
análise de resultados práticos, de forma que o conhecimento seja
construído sequencialmente, como se segue:

\begin{itemize}
\item
  Capítulo 2 - Configurando o Ambiente de Trabalho: Apresenta as
  instruções necessárias para preparar o ambiente de desenvolvimento,
  abordando a integração do \emph{Python} dentro do \emph{RStudio} (via
  pacote \emph{reticulate}) e o gerenciamento de bibliotecas.
\item
  Capítulo 3 - Fundamentos do \emph{Python} para Análise de Dados:
  Revisa as operações básicas, tipos de dados e a criação de funções,
  estabelecendo a base sintática da linguagem.
\item
  Capítulo 4 - Manipulação e Análise Exploratória de Dados: Introduz a
  biblioteca pandas e a criação de visualizações, utilizando conjuntos
  de dados clássicos (como \emph{Titanic} e \emph{Iris}) para demonstrar
  o tratamento e a exploração de informações categóricas e numéricas.
\item
  Capítulo 5 - Computação Numérica com \emph{NumPy}: Foca no uso de
  arrays multidimensionais e operações vetorizadas, ilustrando seu poder
  computacional através de simulações estatísticas.
\item
  Capítulo 6 - Coleta de Dados via \emph{Web Scraping}: Demonstra como
  construir bases de dados próprias automatizando a extração de
  informações de páginas da internet, sempre pautado por boas práticas e
  ética na raspagem de dados.
\item
  Capítulo 7 - Estrutura de Livros Digitais (via \emph{Quarto}): Explica
  brevemente a estrutura de arquivos e códigos para o desenvolvimento de
  livros digitais, utilizando o \emph{Rmarkdown} para os códigos e
  textos e a ferramenta \emph{Quarto} para a compilação.
\end{itemize}

Espera-se que, ao final deste conteúdo, o estudante ou profissional de
Estatística sinta-se confortável para importar, manipular e coletar
dados em \emph{Python}, além de estar apto a construir livros digitais
que podem ser ferramentas ainda mais robustas que os relatórios padrão,
tendo base sólida para estudos posteriores mais aprofundados nestas
ferramentas e suas possibilidades.

\bookmarksetup{startatroot}

\chapter{Configurando o Ambiente de
Trabalho}\label{configurando-o-ambiente-de-trabalho}

Antes de iniciarmos a manipulação de dados e as análises, é fundamental
garantir que o ambiente computacional esteja configurado corretamente.
Para a maioria dos estudantes e profissionais de Estatística, o
\emph{RStudio} é o ambiente de desenvolvimento integrado (IDE) mais
familiar. A boa notícia é que o \emph{RStudio}, especialmente quando
aliado ao \emph{Quarto}, oferece um suporte excelente para a execução de
códigos em \emph{Python}.

\section{\texorpdfstring{O Ambiente \emph{Python} e o pacote
\emph{reticulate}}{O Ambiente Python e o pacote reticulate}}\label{o-ambiente-python-e-o-pacote-reticulate}

Em alguns casos, especialmente na primeira execução de um script
\emph{Python} dentro do \emph{RStudio}, o ambiente pode não estar
perfeitamente configurado. Isso ocorre porque o \emph{RStudio} utiliza
um pacote chamado \emph{reticulate} para fazer a ponte e integrar o
\emph{Python} ao ecossistema do \emph{R.}

O \emph{reticulate} pode acabar criando, de forma automática, um
ambiente virtual \emph{Python} próprio, isolado daquele que você
eventualmente já tenha instalado em seu sistema operacional. Por conta
desse isolamento, esse novo ambiente nem sempre possui todas as
bibliotecas necessárias para a análise de dados já pré-instaladas.

\section{\texorpdfstring{Gerenciamento de Bibliotecas com o
\emph{pip}}{Gerenciamento de Bibliotecas com o pip}}\label{gerenciamento-de-bibliotecas-com-o-pip}

Caso apareçam mensagens de erro indicando que bibliotecas fundamentais
como pandas, \emph{jupyter} ou outras não foram encontradas, elas podem
ser facilmente adicionadas utilizando o \emph{pip.} O \emph{pip} é o
gerenciador de pacotes padrão do \emph{Python}, responsável por baixar,
instalar e atualizar bibliotecas externas.

Para realizar a instalação via \textbf{Terminal} (aba \textbf{Terminal}
no \emph{RStudio}, e não no Console), basta executar os seguintes
comandos:

\begin{Shaded}
\begin{Highlighting}[]
\ExtensionTok{python} \AttributeTok{{-}m}\NormalTok{ pip install pandas numpy seaborn matplotlib}
\ExtensionTok{python} \AttributeTok{{-}m}\NormalTok{ pip install jupyter}
\end{Highlighting}
\end{Shaded}

\section{Instalação Direta em Blocos de
Código}\label{instalauxe7uxe3o-direta-em-blocos-de-cuxf3digo}

Em documentos interativos (como notebooks \emph{Jupyter} ou arquivos
\emph{.qmd} do \emph{Quarto}), também é possível executar comandos de
terminal diretamente de dentro das células de código (\emph{code
chunks}). Para isso, utiliza-se o caractere de exclamação (!) no início
da linha.

Por exemplo, o comando a seguir instala a biblioteca pandas diretamente
a partir do ambiente de execução do documento, sem a necessidade de
abrir o terminal do sistema:

\begin{Shaded}
\begin{Highlighting}[]
\ExtensionTok{!python} \AttributeTok{{-}m}\NormalTok{ pip install pandas}
\end{Highlighting}
\end{Shaded}

\textbf{Dica importante}: Caso o \emph{RStudio} continue apresentando
problemas ou pareça estar utilizando uma instalação do \emph{Python}
diferente da esperada, pode ser necessário reiniciar a sessão do
\emph{R} (\emph{Session} \textgreater{} \emph{Restart R}) para que o
reticulate reconheça as novas bibliotecas instaladas. Em geral, após
essa configuração inicial e a instalação dos pacotes, o código passará a
funcionar normalmente.

\bookmarksetup{startatroot}

\chapter{Fundamentos}\label{fundamentos}

Neste capítulo, iremos contextualizar a lógica de algoritmos ao
\emph{Python}, explicando brevemente a sua sintaxe essencial. A lógica
de programação se mantém a mesma independentemente do ambiente
computacional, mas detalhes estruturais como a tipagem de dados, a
delimitação de blocos e a indexação variam significativamente de uma
linguagem para a outra (Sebesta, 2018).

Enquanto linguagens tipadas estaticamente exigem a declaração prévia do
tipo de cada variável (seja ela um número inteiro ou um texto), o
\emph{Python} utiliza a tipagem dinâmica, inferindo o tipo do dado no
momento de sua execução (Lutz, 2013). Além disso, diferentemente de
linguagens que utilizam chaves para agrupar comandos, o \emph{Python}
adota a indentação obrigatória como regra sintática, o que impõe a
elaboração de códigos padronizados e altamente legíveis. Outro fator de
atenção para os analistas e desenvolvedores é a indexação: ao contrário
do \emph{R} (que é também uma linguagem amplamente utilizada no contexto
de análise de dados e Estatística), onde o primeiro elemento de uma
estrutura de dados ocupa o índice 1, o \emph{Python} adota a indexação
baseada em zero (0-based), característica compartilhada com linguagens
como \emph{C} e o \emph{C++}.

\section{Filosofia, Sintaxe e Tipagem da
Linguagem}\label{filosofia-sintaxe-e-tipagem-da-linguagem}

Grande parte das linguagens de programação surgem (ou evoluem) focadas
em otimização de performance para a máquina. Porém, esse não é o caso do
\emph{Python}, que foi concebido a partir de uma filosofia muito bem
fundamentada e orientada ao desenvolvedor. O princípio norteador da
linguagem é o de que um código é lido muito mais vezes do que ele é
escrito. Portanto, a clareza e a legibilidade devem ser prioridades
absolutas.

Essa filosofia está muito bem documentada no Zen do \emph{Python}
(\texttt{PEP\ 20}), um conjunto de princípios que guiam as melhores
práticas na linguagem, destacando máximas como ``o simples é melhor que
o complexo'' e ``a legibilidade conta'' (Peters, 2004).

Para garantir essa legibilidade de forma estrutural, a sintaxe do
\emph{Python} elimina o uso de marcadores visuais excessivos, como as
chaves \texttt{\{\}} ou as expressões \texttt{begin/end} para delimitar
blocos de código. Em seu lugar, a linguagem adota a indentação
obrigatória. O recuo da margem esquerda não é meramente uma questão de
estilo, mas uma regra sintática rígida documentada na \texttt{PEP\ 8}
(Van Rossum; Warsaw; Coghlan, 2001). Seu objetivo é forçar o programador
a organizar, hierarquicamente, as operações, resultando em um código
mais limpo.

Além da sintaxe visual, o \emph{Python} caracteriza-se por uma tipagem
dinâmica e forte. Dinâmica porque não exige a declaração prévia do tipo
da variável - a própria linguagem infere se o dado é numérico ou textual
no momento da atribuição, além de permitir mudanças de valor e de tipo
durante o processo de desenvolvimento. Forte porque, uma vez inferido o
tipo, o \emph{Python} não realiza conversões implícitas que possam gerar
comportamentos inesperados, garantindo a segurança e a consistência das
operações.

\section{Estruturas de Dados Nativas}\label{estruturas-de-dados-nativas}

Conhecer as estruturas de dados nativas do \emph{Python} é essencial
para armazenar e organizar informações. As principais coleções embutidas
na linguagem são as listas (\texttt{list}), as tuplas (\texttt{tuple}),
os dicionários (\texttt{dict}) e os conjuntos (\texttt{set}).

No contexto da manipulação de dados e da programação geral, a
\textbf{lista} destaca-se como a estrutura mais versátil e
frequentemente utilizada. Diferentemente dos vetores atômicos
encontrados no \emph{R} ou dos \emph{arrays} estáticos comuns no
\emph{C++}, as listas em \emph{Python} são coleções heterogêneas e
dinâmicas. Isso significa que uma única lista pode armazenar, de forma
simultânea, números inteiros, valores reais, textos ou até mesmo outras
listas em seu interior (Lutz, 2013).

Além da heterogeneidade, as listas são estruturas \textbf{mutáveis}.
Após sua criação, o programador ou analista pode alterar seu conteúdo em
tempo de execução sem a necessidade de realocar a estrutura inteira na
memória. É possível modificar valores existentes através de seus
índices, adicionar novos elementos ao final da coleção (utilizando o
método \texttt{append()}) ou remover o último item inserido (através do
método \texttt{pop()}) (Python Software Foundation, 2026a). Essa
flexibilidade de manipulação direta torna as listas uma base poderosa
para o tratamento inicial de dados antes de avançarmos para o uso de
bibliotecas.

\section{Estruturas de Controle de
Fluxo}\label{estruturas-de-controle-de-fluxo}

A espinha dorsal de qualquer linguagem de programação é a capacidade de
direcionar a execução dos algoritmos, através de estruturas de controle
de fluxo. Destacam-se, nesse contexto, as condições lógicas e os laços
de repetição, que são amplamente conhecidos e utilizados por quem
trabalha com códigos, seja qual for o contexto. Em \emph{Python}, as
estruturas de controle de fluxo seguem a filosofia da linguagem - citada
neste capítulo - sendo projetadas para serem diretas e facilmente
legíveis, se comparadas a outras linguagens de programação.

As estruturas condicionais utilizam as palavras-chave \texttt{if},
\texttt{elif} (uma abreviação de \texttt{else\ if}) e \texttt{else}. A
condição avaliada não exige o uso de parênteses, e o bloco de código
subordinado é demarcado de forma simples e intuitiva: através da
inserção de dois pontos (\texttt{:}) ao final da instrução, seguida da
já mencionada indentação obrigatória.

No que diz respeito às estruturas de repetição, o \texttt{Python}
oferece os laços \texttt{while} e \texttt{for}. O laço \texttt{while}
opera de maneira tradicional, executando um bloco de código enquanto uma
condição permanecer verdadeira. Contudo, o laço \texttt{for} apresenta
uma diferença fundamental e vantajosa em relação a linguagens compiladas
clássicas. Em vez de iterar primariamente sobre uma progressão
aritmética de índices, o \texttt{for} em \emph{Python} atua como um
iterador direto sobre os itens de qualquer sequência, como uma lista ou
uma tupla (Python Software Foundation, 2026a).

No âmbito da análise de dados, por exemplo, essa característica permite
que o analista percorra os dados de uma coleção diretamente (por
exemplo, lendo cada observação de uma variável) sem a necessidade de
gerenciar variáveis contadoras, o que reduz significativamente a chance
de erros lógicos. Caso uma iteração numérica tradicional seja
estritamente necessária, a linguagem fornece uma função embutida
\texttt{range()}, capaz de gerar sequências numéricas de forma
eficiente. É importante destacar que a função range() não armazena a
sequência inteira na memória de uma só vez; ela calcula os valores sob
demanda a cada passo da iteração. Consequentemente, sua complexidade de
espaço (consumo de memória) é constante, \(O(1)\)\footnote{A notação
  \(O(1)\) faz parte da chamada ``Notação \emph{Big O}'' (\emph{Big O
  Notation}), utilizada na Ciência da Computação para descrever
  matematicamente a eficiência e a escalabilidade de um algoritmo. Uma
  complexidade de espaço \(O(1)\), ou constante, indica que a quantidade
  de memória exigida para a execução da tarefa permanece exatamente a
  mesma, não importando se a entrada possui dez ou dez milhões de
  elementos.}, independentemente de a sequência possuir dez ou dez
milhões de elementos (Python Software Foundation, 2026b). Isso evita
gargalos de performance e torna o código otimizado.

\section{Funções e Ferramentas
Funcionais}\label{funuxe7uxf5es-e-ferramentas-funcionais}

O processo de desenvolvimento de um algoritmo não apresenta
\textbf{obrigatoriedade} de práticas que vão além de uma lógica concisa
e do respeito à sintaxe da linguagem. Porém, a escalabilidade de um
\emph{software} ou a estrutura de uma análise de dados, por exemplo,
exigem boas práticas que são fundamentadas no bom uso das funções
(métodos), antes de partir para estruturas mais complexas. No
\emph{Python}, a definição de uma função é realizada através da
palavra-chave \texttt{def}, seguida do nome da função e de seus
respectivos parâmetros. Assim como nas funções prontas do
\texttt{Python}, as funções customizadas também exigem, por padrão da
sintaxe, os parênteses logo após o nome da função. Portanto, os
parâmetros são criados dentro dos parênteses, e o valor de saída da
operação, ao final da função, é especificado pelo comando
\texttt{return}.

Contudo, como citado neste capítulo, o \texttt{Python} não se limita ao
paradigma procedural; ele incorpora elementos poderosos do paradigma de
programação funcional. Para evitar a escrita de laços \texttt{for} muito
extensos ao transformar bases de dados, a linguagem fornece funções de
alta ordem (\emph{higher-order functions}), sendo as mais notáveis o
\texttt{map()} e o \texttt{filter()} (Python Software Foundation,
2026c).

Para aqueles que estão habituados ao uso da linguagem \emph{R}, muito
usada na Estatística e na análise de dados, essa mecânica é
conceitualmente idêntica ao uso de funções \texttt{apply()} (como o
\texttt{lapply()} e o \texttt{sapply()}). O objetivo é o mesmo: permitir
um processamento de dados mais rápido, conciso e com menos linhas de
código. O \texttt{map()} aplica uma mesma transformação a todos os itens
de uma coleção de forma simultânea, enquanto o \texttt{filter()} extrai
e retorna apenas os elementos que satisfazem uma determinada condição
lógica (bem próximo do que faz a função \texttt{which()}, no \emph{R}).

Essas ferramentas são ainda mais ágeis quando combinadas com as
expressões \texttt{lambda}. Uma função \texttt{lambda} é uma função
anônima --- ou seja, uma função que não precisa ser formalmente nomeada
com \texttt{def} e que é criada em uma única linha de código, sendo
ideal para operações rápidas ou descartáveis que serão usadas apenas uma
vez dentro de um \texttt{map} ou \texttt{filter} (Lutz, 2013).

\section{Aplicação Prática: Manipulação de
Listas}\label{aplicauxe7uxe3o-pruxe1tica-manipulauxe7uxe3o-de-listas}

No script a seguir, consolidamos todos os fundamentos apresentados neste
capítulo. O exemplo ilustra desde a manipulação direta de uma lista
utilizando métodos nativos (como o \texttt{append()} e o
\texttt{pop()}), até o uso avançado das funções \texttt{map()} e
\texttt{filter()} associadas a funções anônimas (\texttt{lambda}). Por
fim, demonstramos a aplicação de uma função customizada, desenvolvida
com \texttt{def()}, para o cálculo de uma média simples.

Nota-se, também, o uso das chamadas \emph{f-strings}
(\texttt{f"texto\ \{variável\}"}) dentro da função \texttt{print()}.
Esse é o método mais moderno e recomendado no \emph{Python} para
interpretar textos e variáveis na mesma linha, gerando saídas de telas
mais limpas.

\begin{Shaded}
\begin{Highlighting}[]
\CommentTok{\# {-}{-}{-}{-} Definição de uma função customizada {-}{-}{-}{-}}
\KeywordTok{def}\NormalTok{ media(lista):}
\NormalTok{    soma }\OperatorTok{=} \BuiltInTok{sum}\NormalTok{(lista)}
\NormalTok{    quantidade }\OperatorTok{=} \BuiltInTok{len}\NormalTok{(lista)}
    \ControlFlowTok{return}\NormalTok{ soma }\OperatorTok{/}\NormalTok{ quantidade}

\CommentTok{\# {-}{-}{-}{-} Criação e manipulação direta da lista {-}{-}{-}{-}}
\NormalTok{valores }\OperatorTok{=}\NormalTok{ [}\DecValTok{10}\NormalTok{, }\DecValTok{20}\NormalTok{, }\DecValTok{30}\NormalTok{, }\DecValTok{40}\NormalTok{]}
\BuiltInTok{print}\NormalTok{(}\SpecialStringTok{f"Lista original: }\SpecialCharTok{\{}\NormalTok{valores}\SpecialCharTok{\}}\SpecialStringTok{"}\NormalTok{)}

\CommentTok{\# Uso do append()}
\NormalTok{valores.append(}\DecValTok{50}\NormalTok{)}
\BuiltInTok{print}\NormalTok{(}\SpecialStringTok{f"Após append (adicionar 50): }\SpecialCharTok{\{}\NormalTok{valores}\SpecialCharTok{\}}\SpecialStringTok{"}\NormalTok{)}

\CommentTok{\# Uso do pop()}
\NormalTok{valores.pop()}
\BuiltInTok{print}\NormalTok{(}\SpecialStringTok{f"Após pop (remover último elemento): }\SpecialCharTok{\{}\NormalTok{valores}\SpecialCharTok{\}}\SpecialStringTok{"}\NormalTok{)}

\CommentTok{\# {-}{-}{-}{-} Uso de ferramentas funcionais (map, filter e lambda) {-}{-}{-}{-}}

\CommentTok{\# Uso do map() com lambda}
\NormalTok{valores\_dobrados }\OperatorTok{=} \BuiltInTok{list}\NormalTok{(}\BuiltInTok{map}\NormalTok{(}\KeywordTok{lambda}\NormalTok{ x: x }\OperatorTok{*} \DecValTok{2}\NormalTok{, valores))}
\BuiltInTok{print}\NormalTok{(}\SpecialStringTok{f"Valores dobrados com map: }\SpecialCharTok{\{}\NormalTok{valores\_dobrados}\SpecialCharTok{\}}\SpecialStringTok{"}\NormalTok{)}

\CommentTok{\# Uso do filter() com lambda}
\NormalTok{valores\_filtrados }\OperatorTok{=} \BuiltInTok{list}\NormalTok{(}\BuiltInTok{filter}\NormalTok{(}\KeywordTok{lambda}\NormalTok{ x: x }\OperatorTok{\textgreater{}} \DecValTok{20}\NormalTok{, valores))}
\BuiltInTok{print}\NormalTok{(}\SpecialStringTok{f"Valores \textgreater{} 20 com filter: }\SpecialCharTok{\{}\NormalTok{valores\_filtrados}\SpecialCharTok{\}}\SpecialStringTok{"}\NormalTok{)}

\CommentTok{\# {-}{-}{-}{-} Aplicação da função customizada {-}{-}{-}{-}}
\NormalTok{resultado }\OperatorTok{=}\NormalTok{ media(valores)}
\BuiltInTok{print}\NormalTok{(}\SpecialStringTok{f"Média final dos valores originais: }\SpecialCharTok{\{}\NormalTok{resultado}\SpecialCharTok{\}}\SpecialStringTok{"}\NormalTok{)}
\end{Highlighting}
\end{Shaded}

\begin{verbatim}
Lista original: [10, 20, 30, 40]
Após append (adicionar 50): [10, 20, 30, 40, 50]
Após pop (remover último elemento): [10, 20, 30, 40]
Valores dobrados com map: [20, 40, 60, 80]
Valores > 20 com filter: [30, 40]
Média final dos valores originais: 25.0
\end{verbatim}

\bookmarksetup{startatroot}

\chapter{\texorpdfstring{Análise Exploratória e Manipulação de Dados em
\emph{Python}}{Análise Exploratória e Manipulação de Dados em Python}}\label{anuxe1lise-exploratuxf3ria-e-manipulauxe7uxe3o-de-dados-em-python}

\begin{center}\rule{0.5\linewidth}{0.5pt}\end{center}

Chegamos no ponto-chave de qualquer análise de dados. Antes da
modelagem, da inferência e da extração de resultados estatísticos, é a
etapa exploratória que contextualiza o analista. Sem a compreensão
prévia do cenário, qualquer procedimento será puramente matemático:
exato em termos teóricos, porém vazio de informação aplicável ou
interpretável. A função da Estatística é justamente atribuir significado
aos dados, contando a sua história por meio de metodologias robustas e
bem fundamentadas.

Dessa maneira, é preciso compreender a essência das informações em mãos.
É através da Análise Exploratória de Dados (AED) que esse processo ganha
vida e, no ecossistema \emph{Python}, ele é conduzido primordialmente
com o suporte de bibliotecas especializadas na manipulação e
visualização de estruturas tabulares\^{}{[}Uma estrutura tabular
refere-se à organização bidimensional dos dados em formato de tabela. Na
terminologia estatística, equivale a uma matriz de dados, onde cada
linha representa uma observação (ou registro) e cada coluna representa
uma variável (ou atributo) de interesse.

\section{\texorpdfstring{A Biblioteca \emph{Pandas} e as Estruturas de
Dados}{A Biblioteca Pandas e as Estruturas de Dados}}\label{a-biblioteca-pandas-e-as-estruturas-de-dados}

Durante muito tempo, a manipulação de dados tabulares não era uma tarefa
trivial em \emph{Python}, sendo outras linguagens --- como o \emph{R}
--- preferidas para análises estatísticas. Para suprir a necessidade de
ferramentas analíticas flexíveis e de alta performance, o pesquisador
Wes McKinney iniciou, em 2008, o desenvolvimento da biblioteca
\texttt{pandas}. Hoje, ela se consolidou como o padrão ouro da indústria
para a limpeza, transformação e análise de dados em \emph{Python}
(McKinney, 2022). O ecossistema do \texttt{pandas} é construído sobre
duas estruturas de dados fundamentais: a \texttt{Series} e o
\texttt{DataFrame}.

Uma \texttt{Series} pode ser definida como um vetor unidimensional
(semelhante a uma coluna de uma tabela) capaz de armazenar qualquer tipo
de dado, acompanhada por um vetor de rótulos chamado de índice
(\emph{index}).

Já o \texttt{DataFrame} --- a estrutura principal da biblioteca --- é um
objeto bidimensional tabular. Para o estatístico, ele pode ser
compreendido como a equivalência exata de um \texttt{data.frame} na
linguagem \emph{R} ou de uma planilha eletrônica estruturada.
Matematicamente, um \texttt{DataFrame} é essencialmente um dicionário
onde cada coluna é uma \texttt{Series} que compartilha o mesmo índice
posicional (McKinney, 2022).

Apesar da semelhança teórica com o \emph{R}, a manipulação dessas
estruturas apresenta peculiaridades sintáticas importantes. Enquanto no
\emph{R} a indexação se inicia em 1, a do \texttt{pandas} segue a regra
do \emph{Python} (baseada em zero), e a seleção explícita de dados é
intermediada por métodos específicos da biblioteca, como os
localizadores \texttt{.loc} (para rótulos) e \texttt{.iloc} (para
posições inteiras). A partir dessa estrutura bidimensional, o analista
ganha acesso imediato a métodos otimizados para a seleção de colunas,
filtragem de linhas, criação de novas variáveis e tratamento de valores
ausentes.

\section{\texorpdfstring{Importação e Tratamento de Dados (\emph{Data
Wrangling})}{Importação e Tratamento de Dados (Data Wrangling)}}\label{importauxe7uxe3o-e-tratamento-de-dados-data-wrangling}

Antes que qualquer inferência estatística ou visualização seja possível,
os dados brutos precisam ser carregados para a memória e rigorosamente
limpos. Na comunidade analítica, esse processo de estruturação e
higienização é frequentemente chamado de \emph{Data Wrangling}.

O primeiro passo é a importação dos dados. O \texttt{pandas} oferece
diversas opções para a leitura de dados, sendo o método
\texttt{pd.read\_csv()} o mais trivial deles. Ao contrário das funções
básicas de leitura em outras linguagens, o pandas possui um motor de
inferência de tipos altamente otimizado, lidando de forma automática com
a alocação de memória para variáveis numéricas, textuais ou séries
temporais no momento da leitura (VanderPlas, 2016).

Uma vez importada a matriz de dados, é muito comum nos depararmos com
células de valores ausentes (\emph{missing values}). No \emph{R}, estes
valores são representados como \texttt{NA}, enquanto no \emph{Python}
eles são representados como \texttt{NaN}. Ambos representam a mesma
coisa, mudando apenas a forma de apresentação. É evidente que o
tratamento de valores ausentes é um tópico de muita discussão teórica na
Estatística, já que alguns métodos, como o de preenchê-los com a média,
por exemplo, podem nos levar a subestimar a variância em alguns casos, o
que introduziria um viés na análise dos dados. Portanto, este tipo de
tratamento deve sempre ser realizado com atenção considerando que ele
não depende apenas da aplicação de funções computacionais.

Neste contexto, dado que o analista fez, previamente, uma análise
robusta que garanta a viabilidade de algum método de preenchimento de
valores ausentes, o \texttt{pandas} oferece ótimas ferramentas práticas
para a aplicação dos métodos: a função \texttt{dropna()} permite a
exclusão direta das observações ou variáveis incompletas, enquanto a
\texttt{fillna()} viabiliza técnicas de imputação, como o preenchimento
de falhas utilizando a média, a mediana ou o último valor válido
observado da respectiva série.

Outro pilar crítico do Data Wrangling é o cruzamento de bases
relacionais. Projetos reais raramente dependem de um único arquivo
isolado; é frequente a necessidade de unir tabelas distintas por meio de
variáveis-chave, utilizando a função \texttt{pd.merge()}. Nesta fase, a
integridade lógica da base final deve ser a prioridade. Como regra
fundamental de validação de dados em junções, caso ocorram discrepâncias
ou identificadores duplicados ao combinar duas bases, os valores
oriundos da primeira tabela (considerada a base primária ou de maior
confiabilidade) devem sempre prevalecer. No \emph{Python}, esse controle
lógico rigoroso é facilitado pela arquitetura de \emph{joins} do
\texttt{pandas} e por métodos de alinhamento de índices, garantindo que
o tratamento de duplicidades ou conflitos não comprometa a informação
original mais segura.

\section{Visualização dos Dados}\label{visualizauxe7uxe3o-dos-dados}

A qualidade de visualização dos dados é essencial na Estatística. É
através dos recursos gráficos que podemos identificar tendências,
anomalias e até mesmo distribuições. A visualização constitui a fase
inicial de qualquer análise descritiva dos dados, mas não se limita a
apenas isso. É comum, por exemplo, utilizar o QQ-plot para avaliar a
normalidade dos resíduos de um modelo, de forma a dar suporte aos testes
não-paramétricos que podem levar o analista à conclusões equivocadas em
alguns casos, se utilizados de forma isolada. É justamente por este
motivo que a utilização de mais de um recurso é sempre recomendada para
testar pressuposições teóricas importantes.

No âmbito visual, o \emph{Python} conta com bibliotecas muito poderosas.
São elas, principalmente, o \texttt{matplotlib} e o \texttt{Seaborn}. O
\texttt{matplotlib} atua como o motor estrutural da linguagem. Ele é
responsável por instanciar e desenhar os elementos fundamentais da
interface gráfica, como a figura (a ``tela'' em branco), os eixos
cartesianos, as marcações e as escalas (Hunter, 2007). Por ser uma
ferramenta de baixo nível (termo que, na computação, refere-se à menor
distância conceitual entre o código e a máquina, indicando a necessidade
de especificar instruções mais minuciosas, e não uma avaliação sobre a
qualidade da ferramenta) e extremamente detalhista, ele exige uma
quantidade considerável de código para gerar representações
esteticamente polidas, assemelhando-se aos pacotes base de gráficos do R
(como a função \texttt{plot()}).

Para suprir a necessidade de gráficos estatísticos mais expressivos e
com menos linhas de código, surgiu o \texttt{seaborn}. Construída
nativamente sobre o \texttt{matplotlib}, essa biblioteca funciona como
uma interface declarativa de alto nível, sendo perfeitamente integrada
aos objetos da biblioteca \texttt{pandas} (Waskom, 2021).

O grande diferencial do \texttt{seaborn} para o analista de dados é a
abstração de complexidades estruturais somada à automatização de
agregações estatísticas. Em muitas de suas funções, a biblioteca calcula
e exibe intervalos de confiança de forma automática ao agregar dados.
Além disso, o \texttt{seaborn} traz paletas de cores nativas
fundamentadas na teoria da percepção visual, garantindo que as
distinções categóricas e as intensidades contínuas sejam compreendidas
de maneira intuitiva e sem vieses de leitura.

\section{Estudos de Caso Práticos}\label{estudos-de-caso-pruxe1ticos}

Para consolidar os conceitos de \emph{Data Wrangling} e visualização,
aplicaremos as ferramentas do pandas e do \texttt{seaborn} em dois
conjuntos de dados clássicos na literatura estatística e no aprendizado
de máquina: a base \texttt{Titanic} (focada em dados categóricos e taxas
de sobrevivência) e a base \texttt{Iris} (focada em dados contínuos e
distribuições dimensionais).

Para facilitar a reprodutibilidade dos códigos a seguir, utilizaremos as
versões dessas bases que já vêm nativamente embutidas na biblioteca
\texttt{seaborn}.

\subsection{\texorpdfstring{Análise de Dados Categóricos:
\emph{Titanic}}{Análise de Dados Categóricos: Titanic}}\label{anuxe1lise-de-dados-categuxf3ricos-titanic}

O conjunto de dados do naufrágio do \texttt{Titanic} é ideal para a
exploração de tabelas de contingência e distribuições de frequência. No
\emph{Python}, a função \texttt{pd.crosstab()} do pandas atua de forma
muito semelhante à clássica função \texttt{table()} do \emph{R},
permitindo o cruzamento rápido entre variáveis qualitativas.

No script abaixo, importamos as bibliotecas necessárias, carregamos os
dados, removemos rapidamente os valores ausentes de uma coluna de
interesse estrutural e geramos a tabela de contingência cruzando a
classe do passageiro com o status de sobrevivência. Em seguida,
utilizamos o seaborn para criar um gráfico de barras.

\begin{Shaded}
\begin{Highlighting}[]
\ImportTok{import}\NormalTok{ pandas }\ImportTok{as}\NormalTok{ pd}
\ImportTok{import}\NormalTok{ matplotlib.pyplot }\ImportTok{as}\NormalTok{ plt}
\ImportTok{import}\NormalTok{ seaborn }\ImportTok{as}\NormalTok{ sns}

\CommentTok{\# Carregando a base de dados}
\NormalTok{titanic }\OperatorTok{=}\NormalTok{ sns.load\_dataset(}\StringTok{\textquotesingle{}titanic\textquotesingle{}}\NormalTok{)}

\CommentTok{\# Tratamento rápido (Data Wrangling): Removendo linhas sem informação de embarque}
\NormalTok{titanic\_limpo }\OperatorTok{=}\NormalTok{ titanic.dropna(subset}\OperatorTok{=}\NormalTok{[}\StringTok{\textquotesingle{}embark\_town\textquotesingle{}}\NormalTok{])}

\CommentTok{\# Tabela de contingência (Classe vs Sobrevivência)}
\NormalTok{tabela\_sobrevivencia }\OperatorTok{=}\NormalTok{ pd.crosstab(}
\NormalTok{    titanic\_limpo[}\StringTok{\textquotesingle{}class\textquotesingle{}}\NormalTok{], }
\NormalTok{    titanic\_limpo[}\StringTok{\textquotesingle{}survived\textquotesingle{}}\NormalTok{], }
\NormalTok{    rownames}\OperatorTok{=}\NormalTok{[}\StringTok{\textquotesingle{}Classe\textquotesingle{}}\NormalTok{], }
\NormalTok{    colnames}\OperatorTok{=}\NormalTok{[}\StringTok{\textquotesingle{}Sobreviveu?\textquotesingle{}}\NormalTok{]}
\NormalTok{)}
\BuiltInTok{print}\NormalTok{(}\StringTok{"Tabela de Contingência {-} Titanic:}\CharTok{\textbackslash{}n}\StringTok{"}\NormalTok{)}
\BuiltInTok{print}\NormalTok{(tabela\_sobrevivencia)}

\CommentTok{\# Configurando e gerando o gráfico com Seaborn}
\NormalTok{plt.figure(figsize}\OperatorTok{=}\NormalTok{(}\DecValTok{8}\NormalTok{, }\DecValTok{5}\NormalTok{))}
\NormalTok{sns.countplot(}
\NormalTok{    data}\OperatorTok{=}\NormalTok{titanic\_limpo, }
\NormalTok{    x}\OperatorTok{=}\StringTok{\textquotesingle{}class\textquotesingle{}}\NormalTok{, }
\NormalTok{    hue}\OperatorTok{=}\StringTok{\textquotesingle{}survived\textquotesingle{}}\NormalTok{, }
\NormalTok{    palette}\OperatorTok{=}\StringTok{\textquotesingle{}Set2\textquotesingle{}}
\NormalTok{)}
\NormalTok{plt.title(}\StringTok{"Sobrevivência por Classe de Passageiro"}\NormalTok{)}
\NormalTok{plt.xlabel(}\StringTok{"Classe"}\NormalTok{)}
\NormalTok{plt.ylabel(}\StringTok{"Contagem"}\NormalTok{)}
\NormalTok{plt.show()}
\end{Highlighting}
\end{Shaded}

\begin{verbatim}
Tabela de Contingência - Titanic:

Sobreviveu?    0    1
Classe               
First         80  134
Second        97   87
Third        372  119
\end{verbatim}

\pandocbounded{\includegraphics[keepaspectratio]{analise_exploratoria_files/figure-pdf/cell-2-output-2.png}}

\subsection{\texorpdfstring{Análise de Dados Contínuos:
\emph{Iris}}{Análise de Dados Contínuos: Iris}}\label{anuxe1lise-de-dados-contuxednuos-iris}

A base \texttt{Iris}, introduzida pelo estatístico \emph{Ronald Fisher}
em 1936, é o exemplo definitivo para o estudo de distribuições contínuas
e análise de correlação entre múltiplas variáveis numéricas (como
comprimento e largura de pétalas e sépalas).

Antes de partirmos para a visualização, a extração de métricas de
posição e dispersão é o passo primário para a compreensão de variáveis
contínuas. No \emph{R}, essa tarefa é comumente realizada através da
função \texttt{summary()}. No \emph{Python}, a biblioteca
\texttt{pandas} fornece o método \texttt{.describe()}, que calcula
instantaneamente a contagem, a média, o desvio padrão e os quartis de
todas as colunas numéricas do \emph{DataFrame}.

Enquanto as métricas descritivas resumem os dados, a exploração de dados
contínuos exige uma abordagem visual para a identificação de padrões
lineares ou de agrupamentos (\emph{clusters}). O gráfico de dispersão
(\emph{scatterplot}) é a ferramenta analítica primária para esse fim. O
\texttt{seaborn} permite estratificar esses pontos facilmente através do
parâmetro \texttt{hue}, colorindo as observações com base em uma
variável categórica, como a espécie da flor.

\begin{Shaded}
\begin{Highlighting}[]
\CommentTok{\# Carregando a base de dados}
\NormalTok{iris }\OperatorTok{=}\NormalTok{ sns.load\_dataset(}\StringTok{\textquotesingle{}iris\textquotesingle{}}\NormalTok{)}

\CommentTok{\# Extração das métricas descritivas (equivalente ao summary do R)}
\BuiltInTok{print}\NormalTok{(}\StringTok{"Resumo Estatístico da Base Iris:}\CharTok{\textbackslash{}n}\StringTok{"}\NormalTok{)}
\BuiltInTok{print}\NormalTok{(iris.describe())}

\CommentTok{\# Configurando e gerando o gráfico de dispersão com Seaborn}
\NormalTok{plt.figure(figsize}\OperatorTok{=}\NormalTok{(}\DecValTok{8}\NormalTok{, }\DecValTok{5}\NormalTok{))}
\NormalTok{sns.scatterplot(}
\NormalTok{    data}\OperatorTok{=}\NormalTok{iris, }
\NormalTok{    x}\OperatorTok{=}\StringTok{\textquotesingle{}petal\_length\textquotesingle{}}\NormalTok{, }
\NormalTok{    y}\OperatorTok{=}\StringTok{\textquotesingle{}petal\_width\textquotesingle{}}\NormalTok{, }
\NormalTok{    hue}\OperatorTok{=}\StringTok{\textquotesingle{}species\textquotesingle{}}\NormalTok{, }
\NormalTok{    palette}\OperatorTok{=}\StringTok{\textquotesingle{}deep\textquotesingle{}}\NormalTok{,}
\NormalTok{    s}\OperatorTok{=}\DecValTok{80} \CommentTok{\# Tamanho dos pontos}
\NormalTok{)}
\NormalTok{plt.title(}\StringTok{"Relação entre Comprimento e Largura da Pétala"}\NormalTok{)}
\NormalTok{plt.xlabel(}\StringTok{"Comprimento da Pétala (cm)"}\NormalTok{)}
\NormalTok{plt.ylabel(}\StringTok{"Largura da Pétala (cm)"}\NormalTok{)}
\NormalTok{plt.show()}
\end{Highlighting}
\end{Shaded}

\begin{verbatim}
Resumo Estatístico da Base Iris:

       sepal_length  sepal_width  petal_length  petal_width
count    150.000000   150.000000    150.000000   150.000000
mean       5.843333     3.057333      3.758000     1.199333
std        0.828066     0.435866      1.765298     0.762238
min        4.300000     2.000000      1.000000     0.100000
25%        5.100000     2.800000      1.600000     0.300000
50%        5.800000     3.000000      4.350000     1.300000
75%        6.400000     3.300000      5.100000     1.800000
max        7.900000     4.400000      6.900000     2.500000
\end{verbatim}

\pandocbounded{\includegraphics[keepaspectratio]{analise_exploratoria_files/figure-pdf/cell-3-output-2.png}}

\bookmarksetup{startatroot}

\chapter{Referências}\label{referuxeancias}

\phantomsection\label{refs}
\begin{CSLReferences}{0}{1}
\bibitem[\citeproctext]{ref-hunter2007matplotlib}
HUNTER, John D. Matplotlib: A 2D graphics environment. \textbf{Computing
in science \& engineering}, v. 9, n. 3, p. 90--95, 2007.

\bibitem[\citeproctext]{ref-lutz2013learning}
LUTZ, Mark. \textbf{Learning Python: Powerful Object-Oriented
Programming}. 5. ed. \emph{{[}S.l.{]}}: O'Reilly Media, Inc., 2013.

\bibitem[\citeproctext]{ref-mckinney2022python}
MCKINNEY, Wes. \textbf{Python for Data Analysis: Data Wrangling with
pandas, NumPy, and Jupyter}. 3. ed. \emph{{[}S.l.{]}}: O'Reilly Media,
Inc., 2022.

\bibitem[\citeproctext]{ref-peters2004pep20}
PETERS, Tim. \textbf{PEP 20: The Zen of Python}. \emph{{[}S.l.{]}}:
Python Software Foundation, 2004. Disponível em:
\textless{}\url{https://peps.python.org/pep-0020/}\textgreater.

\bibitem[\citeproctext]{ref-python2026tutorial}
PYTHON SOFTWARE FOUNDATION.
\textbf{\href{https://docs.python.org/3/tutorial/datastructures.html}{The
Python Tutorial: Data Structures}}. \emph{{[}S.l.: S.n.{]}}.

\bibitem[\citeproctext]{ref-python2026ranges}
PYTHON SOFTWARE FOUNDATION.
\textbf{\href{https://docs.python.org/3/library/stdtypes.html\#ranges}{Built-in
Types - Ranges}}. \emph{{[}S.l.: S.n.{]}}.

\bibitem[\citeproctext]{ref-python2026builtin}
PYTHON SOFTWARE FOUNDATION.
\textbf{\href{https://docs.python.org/3/library/functions.html}{Built-in
Functions}}. \emph{{[}S.l.: S.n.{]}}.

\bibitem[\citeproctext]{ref-sebesta2018conceitos}
SEBESTA, Robert W. \textbf{Conceitos de Linguagens de Programação}. 11.
ed. \emph{{[}S.l.{]}}: Bookman Editora, 2018.

\bibitem[\citeproctext]{ref-vanrossum2001pep8}
VAN ROSSUM, Guido; WARSAW, Barry; COGHLAN, Nick. \textbf{PEP 8: Style
Guide for Python Code}. \emph{{[}S.l.{]}}: Python Software Foundation,
2001. Disponível em:
\textless{}\url{https://peps.python.org/pep-0008/}\textgreater.

\bibitem[\citeproctext]{ref-vanderplas2016python}
VANDERPLAS, Jake. \textbf{Python Data Science Handbook: Essential Tools
for Working with Data}. \emph{{[}S.l.{]}}: O'Reilly Media, Inc., 2016.

\bibitem[\citeproctext]{ref-waskom2021seaborn}
WASKOM, Michael L. seaborn: statistical data visualization.
\textbf{Journal of Open Source Software}, v. 6, n. 60, p. 3021, 2021.

\end{CSLReferences}




\end{document}
